\section{总结与展望}\label{sec:conclusion}

\subsection{工作总结}

本文围绕遥感图像中的微小船舶检测，完成了文献梳理、模型实验和系统实现几项工作。研究部分先整理通用目标检测、小目标检测和遥感船舶检测的发展历史，再在同一数据和评测口径下比较 DRENet、FCOS 与 YOLO26 模型。工程部分采用 React、FastAPI 和 SQLite 搭建图像检测系统，完成任务提交、异步推理、结果可视化和历史回看等功能。

实验部分的主要结论是，DRENet 和 YOLO26 模型的 AP50 非常接近，但二者取向不同。DRENet 的 Recall 更高，适合需要优先发现目标的场景；YOLO26 在 AP50-95、Precision 和 F1 上则更加均衡。而 FCOS 的整体结果弱于前两者，但保留了 anchor-free 这一路线的参照意义。阈值消融和输入尺寸敏感性分析也充分说明，误检和漏检不能只靠一个 AP50 指标解释。

系统部分并不是单独做一个结果展示页，而是把推理能力、任务状态、结果文件和历史记录放到同一套流程里管理。经过联调后，系统已经可以接入真实模型权重，支持任务回放，离线实验中的部分结果也能在页面中复查。

本文的工作重点不在提出新的检测网络，而在于把数据口径、实验记录和系统实现整理成一条可复查的路径。

\subsection{不足分析}

当前工作已经完成主要目标，但仍有几处边界。

第一，DRENet、FCOS 与 YOLO26 的不同实现路线仍有进一步统一空间。第二，融合模式已经跑通，但还可以继续补充更系统的离线评测。第三，系统层面仍可以继续完善部署说明和跨环境迁移流程。

这些问题更多是后续可以继续补充的方向，不影响本文对 DRENet、YOLO26 和 FCOS 的主要结论。如果要继续提高可迁移性和评测严谨性，重点应该放在部署规范、统一离线评估和跨模型对齐上。

\subsection{未来展望}

后续研究可从以下方面展开：把融合模式的统一离线评测补齐，把单模型结果和融合展示之间的关系量化出来；继续围绕 DRENet 和 YOLO26 的误检、漏检取舍来调阈值、后处理和融合策略；继续整理 FCOS 的训练配置，让跨模型比较更严谨；完善模型接入和部署流程；在更多遥感场景和数据集上检查方法和系统的泛化情况。

沿着这些方向继续做，现有系统可以进一步扩展成实验管理平台。比如把统一报告导出、批量回评和参数留痕都接到同一个工作流里，让实验复查和系统展示共享更多底层能力。

\subsection{全文结论}

本文完成了遥感微小船舶检测任务中的实验比较和系统实现。在 LEVIR-Ship 数据集上，DRENet 更偏召回优先，YOLO26 更偏精度和部署均衡，FCOS 保留为 anchor-free 参照。系统实现则说明，将推理能力服务化并纳入任务管理后，结果展示和历史回放会更加清楚。

从最终完成情况看，本文已经把问题分析、实验论证和系统实现串联起来。现有实验记录、分析框架和系统代码，为后续性能优化、评测补齐和部署规范化留下了可继续迭代的基础。
